% Compilar dos veces con pdflatex desde este directorio.
\documentclass[aspectratio=169,10pt]{beamer}
\usepackage[utf8]{inputenc}
\usepackage[T1]{fontenc}
\usepackage[spanish,es-nodecimaldot]{babel}
\usepackage{amsmath,amssymb,booktabs,tabularx,array,tikz}
\usetikzlibrary{arrows.meta,positioning,calc,shapes.geometric,fit}

\definecolor{UAPNavy}{HTML}{17324D}
\definecolor{UAPTeal}{HTML}{007F82}
\definecolor{UAPGold}{HTML}{E2A33A}
\definecolor{UAPLight}{HTML}{F3F6F8}
\definecolor{UAPGray}{HTML}{536471}
\definecolor{UAPRed}{HTML}{A94343}

\tikzset{
  flow/.style={-{Stealth[length=2.2mm]},thick,draw=UAPNavy},
  feedback/.style={{Stealth[length=2.2mm]}-{Stealth[length=2.2mm]},thick,draw=UAPNavy},
  datum/.style={draw=UAPGold,fill=UAPGold!16,rounded corners=2pt,align=center,minimum height=8mm,text=UAPNavy},
  process/.style={draw=UAPTeal,fill=UAPTeal!12,rounded corners=3pt,very thick,align=center,minimum height=10mm,text=UAPNavy},
  decision/.style={draw=UAPNavy,fill=UAPNavy!9,rounded corners=3pt,very thick,align=center,minimum height=10mm,text=UAPNavy},
  warning/.style={draw=UAPRed,fill=UAPRed!7,rounded corners=3pt,align=left},
  neuron/.style={circle,draw=UAPTeal,fill=UAPTeal!12,minimum size=7mm,inner sep=0pt},
  client/.style={draw=UAPGold,fill=UAPGold!14,rounded corners=2pt,align=center,minimum width=24mm,minimum height=9mm}
}

\usetheme{Boadilla}
\usecolortheme{default}
\setbeamercolor{normal text}{fg=UAPNavy,bg=white}
\setbeamercolor{structure}{fg=UAPTeal}
\setbeamercolor{frametitle}{fg=white,bg=UAPNavy}
\setbeamercolor{title}{fg=white}
\setbeamercolor{subtitle}{fg=UAPGold}
\setbeamercolor{block title}{fg=white,bg=UAPTeal}
\setbeamercolor{block body}{fg=UAPNavy,bg=UAPLight}
\setbeamercolor{alertblock title}{fg=white,bg=UAPRed}
\setbeamercolor{alertblock body}{fg=UAPNavy,bg=UAPRed!7}
\setbeamerfont{frametitle}{series=\bfseries}
\setbeamertemplate{navigation symbols}{}
\setbeamertemplate{itemize item}{\color{UAPGold}$\blacktriangleright$}
\setbeamertemplate{itemize subitem}{\color{UAPTeal}$\bullet$}
\setbeamertemplate{footline}{%
  \leavevmode\hbox{%
    \begin{beamercolorbox}[wd=.78\paperwidth,ht=2.8ex,dp=1.2ex,leftskip=1em]{author in head/foot}%
      \color{white}\insertshorttitle
    \end{beamercolorbox}%
    \begin{beamercolorbox}[wd=.22\paperwidth,ht=2.8ex,dp=1.2ex,rightskip=1em plus1fil]{date in head/foot}%
      \color{UAPGray}\hfill Semana 8 \enspace | \enspace \insertframenumber/\inserttotalframenumber
    \end{beamercolorbox}}}

\newcommand{\concepto}[1]{\textcolor{UAPTeal}{\textbf{#1}}}
\newcommand{\idea}[1]{\begin{block}{Idea central}#1\end{block}}
\title[IA · Clase 8]{Panorama de métodos de Inteligencia Artificial}
\subtitle{Modelos, interacción, lenguaje y aprendizaje distribuido}
\author{Curso de Inteligencia Artificial}
\institute{Semana 8 · Clase 8}
\date{}

\begin{document}

% 1
\begin{frame}[plain]
  \begin{beamercolorbox}[wd=\paperwidth,ht=\paperheight,sep=2em]{frametitle}
    \vfill
    {\Huge\bfseries\inserttitle\par}\vspace{.45em}
    {\color{UAPGold}\Large\insertsubtitle\par}\vspace{1.2em}
    \begin{tikzpicture}[node distance=4mm,scale=.69,transform shape]
      \node[datum] (s) {Señal};
      \node[process,right=of s] (r) {Representación};
      \node[datum,right=of r] (a) {Aprendizaje};
      \node[decision,right=of a] (o) {Salida};
      \node[datum,right=of o] (e) {Evaluación};
      \draw[flow] (s)--(r); \draw[flow] (r)--(a); \draw[flow] (a)--(o); \draw[flow] (o)--(e);
    \end{tikzpicture}\par
    \vfill {\color{UAPGold}\large\insertauthor\par}{\color{white}\insertinstitute\par}
  \end{beamercolorbox}
\end{frame}

% 2
\begin{frame}{Propósito de la clase}
  Comprender los principios que distinguen varias familias centrales de Inteligencia Artificial y reconocer qué aprende cada una.
  \begin{block}{Al finalizar, podremos}
    \begin{itemize}
      \item explicar SVM, redes neuronales y ensambles;
      \item describir el ciclo de aprendizaje por refuerzo;
      \item seguir el recorrido desde texto hasta un modelo de lenguaje;
      \item explicar el aprendizaje federado;
      \item comparar señales, salidas, fortalezas y límites.
    \end{itemize}
  \end{block}
\end{frame}

% 3
\begin{frame}{SVM: separar con margen}
  Una máquina de soporte vectorial busca una frontera que separe clases y deje el mayor margen posible.
  \[
    f(\mathbf{x})=\mathbf{w}^{\mathsf T}\mathbf{x}+b,
    \qquad \widehat y=\operatorname{sign}(f(\mathbf{x})).
  \]
  \begin{itemize}
    \item La frontera está dada por $f(\mathbf{x})=0$.
    \item $\mathbf{w}$ determina su orientación y $b$ su desplazamiento.
    \item La distancia firmada es proporcional a $f(\mathbf{x})/\|\mathbf{w}\|_2$.
  \end{itemize}
\end{frame}

% 6
\begin{frame}{Margen y vectores de soporte}
  \begin{columns}[T]
    \column{.52\textwidth}
    \centering
    \begin{tikzpicture}[scale=.72]
      \draw[->,UAPGray] (0,0)--(6.4,0) node[right] {$x_1$};
      \draw[->,UAPGray] (0,0)--(0,4.8) node[above] {$x_2$};
      \draw[UAPGold,dashed,very thick] (.7,3.8)--(4.7,.2);
      \draw[UAPNavy,very thick] (1.4,4.6)--(5.4,.6) node[pos=.57,above right] {frontera};
      \draw[UAPGold,dashed,very thick] (2.3,4.8)--(6.1,1.0);
      \foreach \x/\y in {.7/1.0,1.2/1.7,1.9/.8,2.0/2.0,2.8/1.5}{\fill[UAPTeal] (\x,\y) circle (2.3pt);}
      \foreach \x/\y in {3.2/4.1,4.0/3.5,4.8/4.2,5.3/3.1,5.7/4.4}{\fill[UAPRed] (\x,\y) circle (2.3pt);}
      \draw[UAPNavy,thick] (2.0,2.0) circle (4.4pt);
      \draw[UAPNavy,thick] (3.2,4.1) circle (4.4pt);
      \draw[<->,UAPGray] (2.45,2.52)--(3.0,3.08) node[midway,right] {margen};
    \end{tikzpicture}
    \column{.46\textwidth}
    \small
    \begin{itemize}
      \item Ancho del margen: $2/\|\mathbf{w}\|_2$.
      \item Los ejemplos marcados son \concepto{vectores de soporte}.
      \item Determinan la posición de la frontera.
      \item Casos alejados suelen no modificar la solución.
    \end{itemize}
  \end{columns}
\end{frame}

% 7
\begin{frame}{Margen blando y kernels}
  \small
  En datos reales puede ser imposible separar sin errores. El parámetro $C$ regula el equilibrio entre margen amplio y penalización de violaciones.
  \[
    K(\mathbf{x},\mathbf{z})=\phi(\mathbf{x})^{\mathsf T}\phi(\mathbf{z}).
  \]
  \begin{tabularx}{\textwidth}{>{\bfseries}p{.25\textwidth}XX}\toprule
  Opción & Qué representa & Cuidado\\\midrule
  Kernel lineal & frontera lineal disponible & depende de escala y variables\\
  Kernel RBF & semejanza local y frontera curva & requiere elegir $C$ y $\gamma$\\
  Margen blando & admite violaciones & no convierte ruido en señal\\\bottomrule
  \end{tabularx}
\end{frame}

% 8
\begin{frame}{Ejemplo SVM: riesgo operativo}
  \small
  Se representan eventos mediante demora estimada y congestión.
  \begin{enumerate}
    \item Un SVM lineal busca una recta entre riesgo bajo y alto.
    \item Los eventos cercanos son vectores de soporte.
    \item Si la separación es curva, un kernel RBF puede representarla.
    \item La puntuación respecto de la frontera no es automáticamente una probabilidad calibrada.
  \end{enumerate}
  \begin{alertblock}{Límite}La acción posterior todavía requiere costos, capacidad y supervisión.\end{alertblock}
\end{frame}

% 9
\begin{frame}{Red neuronal: una unidad básica}
  \begin{columns}[T]
    \column{.54\textwidth}
    Una neurona combina entradas y aplica una activación:
    \[
      z=\mathbf{w}^{\mathsf T}\mathbf{x}+b,
      \qquad a=\phi(z).
    \]
    \begin{itemize}
      \item Los pesos controlan la influencia de cada entrada.
      \item El sesgo desplaza el umbral.
      \item La activación introduce una transformación.
    \end{itemize}
    \column{.42\textwidth}
    \centering
    \begin{tikzpicture}[node distance=6mm]
      \node[datum] (x1) {$x_1$};
      \node[datum,below=of x1] (x2) {$x_2$};
      \node[datum,below=of x2] (xp) {$x_p$};
      \node[process,right=15mm of x2] (sum) {$\sum_j w_jx_j+b$};
      \node[decision,right=8mm of sum] (act) {$\phi(z)$};
      \draw[flow] (x1)--(sum); \draw[flow] (x2)--(sum); \draw[flow] (xp)--(sum); \draw[flow] (sum)--(act);
    \end{tikzpicture}
  \end{columns}
\end{frame}

% 10
\begin{frame}{Funciones de activación}
  \small
  \begin{tabularx}{\textwidth}{>{\bfseries}p{.22\textwidth}p{.29\textwidth}X}\toprule
  Activación & Expresión o rango & Uso frecuente\\\midrule
  Sigmoide & $\sigma(z)=1/(1+e^{-z})$ & salida binaria\\
  Tanh & $(-1,1)$ & estados centrados\\
  ReLU & $\max(0,z)$ & capas ocultas\\
  Softmax & probabilidades que suman 1 & salida multiclase\\\bottomrule
  \end{tabularx}
  \begin{block}{Por qué importan}Sin activaciones no lineales, apilar capas seguiría produciendo una transformación lineal.\end{block}
\end{frame}

% 11
\begin{frame}{Capas y propagación hacia adelante}
  \begin{columns}[T]
    \column{.49\textwidth}
    \centering
    \begin{tikzpicture}[x=1.4cm,y=.8cm]
      \foreach \i in {1,...,3}{\node[neuron] (i\i) at (0,2-\i) {};}
      \foreach \i in {1,...,4}{\node[neuron] (h\i) at (2,2.5-\i) {};}
      \foreach \i in {1,...,2}{\node[neuron] (o\i) at (4,1.5-\i) {};}
      \foreach \i in {1,...,3}\foreach \j in {1,...,4}{\draw[UAPGray!55] (i\i)--(h\j);}
      \foreach \i in {1,...,4}\foreach \j in {1,...,2}{\draw[UAPGray!55] (h\i)--(o\j);}
      \node[below=5mm of i3] {entrada}; \node[below=1mm of h4] {oculta}; \node[below=5mm of o2] {salida};
      \draw[flow,UAPGold,very thick] (-.3,2)--(4.3,2) node[midway,above] {propagación};
    \end{tikzpicture}
    \column{.49\textwidth}
    \small
    \[
      \mathbf{h}^{(\ell)}=\phi\!\left(W^{(\ell)}\mathbf{h}^{(\ell-1)}+\mathbf{b}^{(\ell)}\right).
    \]
    \begin{itemize}
      \item Entrada: variables observadas.
      \item Ocultas: representaciones aprendidas.
      \item Salida: cantidad, clase o probabilidad.
    \end{itemize}
    La inferencia calcula una salida; no modifica parámetros.
  \end{columns}
\end{frame}

% 12
\begin{frame}{Pérdida y retropropagación}
  \begin{center}
    \begin{tikzpicture}[node distance=5mm,scale=.78,transform shape]
      \node[datum] (x) {Entrada};
      \node[process,right=of x] (f) {Forward};
      \node[decision,right=of f] (y) {Predicción};
      \node[warning,right=of y] (l) {Pérdida};
      \node[process,below=8mm of y] (b) {Gradientes\\backprop};
      \node[datum,left=of b] (u) {Actualizar\\parámetros};
      \draw[flow] (x)--(f); \draw[flow] (f)--(y); \draw[flow] (y)--(l);
      \draw[flow] (l)--(b); \draw[flow] (b)--(u); \draw[flow] (u)--(f);
    \end{tikzpicture}
  \end{center}
  \[
    \theta\leftarrow\theta-\eta\nabla_{\theta}L.
  \]
  \small La retropropagación calcula gradientes; el optimizador los usa para ajustar parámetros. La validación orienta arquitectura y parada; la prueba estima desempeño final.
\end{frame}

% 13
\begin{frame}{Ejemplo neuronal: relación XOR}
  \begin{columns}[T]
    \column{.38\textwidth}
    \centering
    \begin{tabular}{ccc}\toprule
    $x_1$ & $x_2$ & $y$\\\midrule
    0&0&0\\0&1&1\\1&0&1\\1&1&0\\\bottomrule
    \end{tabular}
    \column{.58\textwidth}
    \begin{itemize}
      \item Una sola recta no separa los casos.
      \item Una capa oculta construye regiones intermedias.
      \item La salida combina esas regiones en una frontera no lineal.
    \end{itemize}
    \begin{block}{Resultado}Capacidad para ajustar no equivale a generalización: una red también puede memorizar ruido.\end{block}
  \end{columns}
\end{frame}

% 14
\begin{frame}{Ensemble Learning: combinar modelos}
  Un ensamble agrega predicciones de varios modelos:
  \[
    \widehat f(\mathbf{x})=\sum_{m=1}^{M}\alpha_m f_m(\mathbf{x}).
  \]
  \begin{itemize}
    \item La combinación puede ser promedio, voto, suma ponderada o metamodelo.
    \item La mejora aparece cuando los modelos son competentes y sus errores difieren.
    \item Combinar copias idénticas con los mismos errores aporta poco.
  \end{itemize}
\end{frame}

% 15
\begin{frame}{Bagging y Random Forest}
  \begin{center}
    \begin{tikzpicture}[node distance=5mm,scale=.8,transform shape]
      \node[datum] (d) {Datos};
      \node[process,right=10mm of d,yshift=12mm] (m1) {Muestra 1\\Árbol 1};
      \node[process,right=10mm of d] (m2) {Muestra 2\\Árbol 2};
      \node[process,right=10mm of d,yshift=-12mm] (m3) {Muestra 3\\Árbol 3};
      \node[decision,right=14mm of m2] (a) {Promedio\\o voto};
      \draw[flow] (d)--(m1); \draw[flow] (d)--(m2); \draw[flow] (d)--(m3);
      \draw[flow] (m1)--(a); \draw[flow] (m2)--(a); \draw[flow] (m3)--(a);
    \end{tikzpicture}
  \end{center}
  \small Bagging entrena en paralelo sobre muestras bootstrap. Random Forest añade selección aleatoria de variables para reducir la correlación entre árboles.
\end{frame}

% 16
\begin{frame}{Boosting}
  \begin{center}
    \begin{tikzpicture}[node distance=6mm,scale=.82,transform shape]
      \node[process] (m1) {Modelo 1};
      \node[warning,right=of m1] (e1) {Error 1};
      \node[process,right=of e1] (m2) {Modelo 2};
      \node[warning,right=of m2] (e2) {Error 2};
      \node[process,right=of e2] (m3) {Modelo 3};
      \draw[flow] (m1)--(e1); \draw[flow] (e1)--(m2); \draw[flow] (m2)--(e2); \draw[flow] (e2)--(m3);
    \end{tikzpicture}
  \end{center}
  \[
    F_M(x)=F_0(x)+\sum_{m=1}^{M}\nu h_m(x).
  \]
  Cada modelo corrige parte del error actual. No es una votación independiente: el aprendizaje es secuencial.
\end{frame}

% 17
\begin{frame}{Ejemplo de ensamble}
  \small
  Para estimar riesgo de demora:
  \begin{tabularx}{\textwidth}{>{\bfseries}p{.24\textwidth}XX}\toprule
  Método & Mecanismo & Efecto buscado\\\midrule
  Árbol & reglas por regiones & modelo base interpretable\\
  Random Forest & árboles decorrelacionados & menor variabilidad\\
  Boosting & correcciones sucesivas & menor error sistemático\\
  Stacking & metamodelo sobre predicciones & combinación aprendida\\\bottomrule
  \end{tabularx}
  \idea{La comparación conserva la misma partición, información y métrica.}
\end{frame}

% 18
\begin{frame}{Aprendizaje por refuerzo}
  En aprendizaje supervisado, una etiqueta indica la respuesta esperada. En RL, un agente actúa, observa consecuencias y aprende una política.
  \[
    S_t\rightarrow A_t\rightarrow R_{t+1},S_{t+1}.
  \]
  \begin{itemize}
    \item La señal puede llegar con demora.
    \item Una acción cambia las experiencias futuras.
    \item El objetivo es una secuencia de decisiones, no una predicción aislada.
  \end{itemize}
\end{frame}

% 19
\begin{frame}{Ciclo agente-entorno}
  \begin{center}
    \begin{tikzpicture}[node distance=28mm]
      \node[process,minimum width=35mm,minimum height=18mm] (agent) {Agente\\política / valor};
      \node[decision,right=of agent,minimum width=35mm,minimum height=18mm] (env) {Entorno\\transición};
      \draw[flow,bend left=18] (agent) to node[above] {acción $A_t$} (env);
      \draw[flow,bend left=18] (env) to node[below] {estado $S_{t+1}$, recompensa $R_{t+1}$} (agent);
    \end{tikzpicture}
  \end{center}
  \begin{itemize}
    \item El estado resume información relevante.
    \item La recompensa aporta una señal inmediata.
    \item La política asigna acciones a estados.
  \end{itemize}
\end{frame}

% 20
\begin{frame}{Retorno y funciones de valor}
  \small
  \[
    G_t=R_{t+1}+\gamma R_{t+2}+\gamma^2R_{t+3}+\cdots,
    \qquad 0\leq\gamma<1.
  \]
  \begin{columns}[T]
    \column{.48\textwidth}
    \begin{block}{Valor de estado}
      $V^{\pi}(s)=\mathbb{E}_{\pi}[G_t\mid S_t=s]$
    \end{block}
    \column{.48\textwidth}
    \begin{block}{Valor de acción}
      $Q^{\pi}(s,a)=\mathbb{E}_{\pi}[G_t\mid S_t=s,A_t=a]$
    \end{block}
  \end{columns}
  $\gamma$ regula cuánto pesan consecuencias futuras; el valor resume retorno esperado.
\end{frame}

% 21
\begin{frame}{Exploración y explotación}
  \begin{columns}[T]
    \column{.48\textwidth}
    \begin{block}{Explotar}
      Elegir la acción que actualmente parece mejor. Puede fijar una política prematura.
    \end{block}
    \column{.48\textwidth}
    \begin{block}{Explorar}
      Probar alternativas para obtener información. Puede ser costoso o inseguro.
    \end{block}
  \end{columns}
  \vspace{4mm}
  La estrategia $\varepsilon$-greedy elige una acción aleatoria con probabilidad $\varepsilon$ y la mejor conocida en los demás casos.
\end{frame}

% 22
\begin{frame}{Q-learning como ejemplo}
  \[
    Q(s,a)\leftarrow Q(s,a)+\alpha\left[r+\gamma\max_{a'}Q(s',a')-Q(s,a)\right].
  \]
  \begin{center}
    \begin{tikzpicture}[node distance=5mm,scale=.85,transform shape]
      \node[datum] (old) {$Q(s,a)$ actual};
      \node[process,right=of old] (target) {$r+\gamma\max Q(s',a')$};
      \node[decision,right=of target] (new) {$Q(s,a)$ actualizado};
      \draw[flow] (old)--node[above] {comparar} (target); \draw[flow] (target)--node[above] {$\alpha$} (new);
    \end{tikzpicture}
  \end{center}
  El objetivo combina recompensa inmediata y mejor valor futuro estimado.
\end{frame}

% 23
\begin{frame}{Ejemplo RL: asignación secuencial}
  \small
  \begin{tabularx}{\textwidth}{>{\bfseries}p{.22\textwidth}X}\toprule
  Elemento & Ejemplo\\\midrule
  Estado & zona, hora, congestión y disponibilidad\\
  Acción & mantener, mover o esperar\\
  Recompensa & servicio logrado menos traslado y demora\\
  Política & regla aprendida para seleccionar acciones\\\bottomrule
  \end{tabularx}
  \begin{alertblock}{Diseño de la señal}Una recompensa incompleta puede producir conductas indeseadas aunque el algoritmo la optimice correctamente.\end{alertblock}
\end{frame}

% 24
\begin{frame}{NLP: del texto a una tarea}
  \begin{center}
    \begin{tikzpicture}[node distance=5mm,scale=.78,transform shape]
      \node[datum] (t) {Texto};
      \node[process,right=of t] (tok) {Tokens};
      \node[process,right=of tok] (rep) {Representación\\numérica};
      \node[datum,right=of rep] (mod) {Modelo};
      \node[decision,right=of mod] (out) {Clase, entidad,\\respuesta o texto};
      \draw[flow] (t)--(tok); \draw[flow] (tok)--(rep); \draw[flow] (rep)--(mod); \draw[flow] (mod)--(out);
    \end{tikzpicture}
  \end{center}
  \small Un corpus es una colección de documentos; un token es una unidad procesada. Los tokens pueden ser palabras, subpalabras, caracteres o símbolos especiales.
\end{frame}

% 25
\begin{frame}{Representaciones dispersas}
  Bag of Words representa documentos mediante conteos. TF--IDF reduce el peso de términos frecuentes en muchos documentos:
  \[
    \operatorname{tfidf}(t,d)=\operatorname{tf}(t,d)\log\frac{N}{\operatorname{df}(t)}.
  \]
  \begin{itemize}
    \item Son baselines útiles e interpretables.
    \item Producen vectores grandes y dispersos.
    \item Pierden gran parte del orden y del contexto.
  \end{itemize}
\end{frame}

% 26
\begin{frame}{Embeddings}
  \[
    \text{token}\longrightarrow\mathbf{e}\in\mathbb{R}^{d}.
  \]
  \begin{itemize}
    \item Un embedding representa una unidad mediante un vector denso.
    \item Los embeddings distribucionales aproximan semejanzas de uso.
    \item Los embeddings contextuales cambian según la oración.
  \end{itemize}
  \begin{block}{Ejemplo}“Banco” necesita representaciones distintas en “banco financiero” y “banco de la plaza”.\end{block}
\end{frame}

% 27
\begin{frame}{Atención}
  \[
    \operatorname{Attention}(Q,K,V)=
    \operatorname{softmax}\left(\frac{QK^{\mathsf T}}{\sqrt{d_k}}\right)V.
  \]
  \begin{center}
    \begin{tikzpicture}[node distance=7mm,scale=.82,transform shape]
      \node[datum] (q) {Consulta $Q$\\qué busca};
      \node[datum,right=of q] (k) {Claves $K$\\con qué compara};
      \node[process,right=of k] (w) {Pesos\\softmax};
      \node[datum,right=of w] (v) {Valores $V$\\qué combina};
      \node[decision,right=of v] (o) {Contexto};
      \draw[flow] (q)--(w); \draw[flow] (k)--(w); \draw[flow] (w)--(v); \draw[flow] (v)--(o);
    \end{tikzpicture}
  \end{center}
\end{frame}

% 28
\begin{frame}{Transformer simplificado}
  \begin{center}
    \begin{tikzpicture}[node distance=4mm,scale=.76,transform shape]
      \node[datum] (tok) {Tokens};
      \node[process,right=of tok] (emb) {Embeddings\\+ posición};
      \node[process,right=of emb] (att) {Atención\\multi-cabeza};
      \node[process,right=of att] (ff) {Red\\feed-forward};
      \node[datum,right=of ff] (stack) {Capas\\apiladas};
      \node[decision,right=of stack] (out) {Salida};
      \draw[flow] (tok)--(emb); \draw[flow] (emb)--(att); \draw[flow] (att)--(ff); \draw[flow] (ff)--(stack); \draw[flow] (stack)--(out);
    \end{tikzpicture}
  \end{center}
  \begin{itemize}
    \item Conexiones residuales y normalización estabilizan el procesamiento.
    \item La atención permite intercambiar información entre posiciones.
    \item El diagrama omite detalles, pero conserva el flujo funcional.
  \end{itemize}
\end{frame}

% 29
\begin{frame}{Modelos de lenguaje}
  Un modelo de lenguaje asigna probabilidades a secuencias. En generación autorregresiva:
  \[
    P(x_{1:T})=\prod_{t=1}^{T}P(x_t\mid x_{<t}).
  \]
  \begin{itemize}
    \item Un LLM amplía escala de parámetros, datos y cómputo.
    \item El modelo produce una distribución sobre tokens posibles.
    \item Una aplicación conversacional añade interfaz, instrucciones, memoria y herramientas.
  \end{itemize}
\end{frame}

% 30
\begin{frame}{Ejemplo de generación token a token}
  \begin{columns}[T]
    \column{.52\textwidth}
    Contexto: \concepto{“La capital de Argentina es”}
    \begin{center}
      \begin{tabular}{lr}\toprule
      Token candidato & Probabilidad\\\midrule
      Buenos & 0,72\\
      Córdoba & 0,08\\
      una & 0,05\\
      grande & 0,02\\\bottomrule
      \end{tabular}
    \end{center}
    \column{.44\textwidth}
    \begin{tikzpicture}[node distance=5mm]
      \node[datum] (c) {Contexto};
      \node[process,below=of c] (m) {Modelo};
      \node[decision,below=of m] (p) {Distribución};
      \node[datum,below=of p] (n) {Nuevo token};
      \draw[flow] (c)--(m); \draw[flow] (m)--(p); \draw[flow] (p)--(n);
      \draw[flow] (n.east) -- ++(.7,0) |- (c.east);
    \end{tikzpicture}
  \end{columns}
\end{frame}

% 31
\begin{frame}{Fluidez no equivale a verdad}
  \begin{alertblock}{Límite del objetivo}
    Predecir tokens probables no verifica automáticamente afirmaciones.
  \end{alertblock}
  \begin{itemize}
    \item Una salida probable puede ser falsa.
    \item El contexto puede ser insuficiente o ambiguo.
    \item Los datos pueden contener sesgos o información desactualizada.
    \item Recuperar fuentes ayuda, pero no elimina la validación.
  \end{itemize}
  La evaluación debe corresponder al uso: exactitud factual, cobertura, robustez, seguridad y trazabilidad.
\end{frame}

% 32
\begin{frame}{Aprendizaje federado}
  El aprendizaje federado entrena un modelo compartido sin reunir todos los datos brutos en un único repositorio.
  \begin{center}
    \concepto{modelo global $\rightarrow$ entrenamiento local $\rightarrow$ actualizaciones $\rightarrow$ agregación}
  \end{center}
  \begin{itemize}
    \item Es una organización distribuida del entrenamiento.
    \item No es una familia de predicción.
    \item El modelo local puede ser una red u otra arquitectura compatible.
  \end{itemize}
\end{frame}

% 33
\begin{frame}{Una ronda de FedAvg}
  \begin{columns}[T]
    \column{.54\textwidth}
    \centering
    \begin{tikzpicture}[node distance=7mm,scale=.72,transform shape]
      \node[process,minimum width=30mm] (server) {Servidor\\$\theta_t\rightarrow\theta_{t+1}$};
      \node[client,below left=14mm and 7mm of server] (c1) {Cliente 1\\$n_1,\theta^{(1)}$};
      \node[client,below=14mm of server] (c2) {Cliente 2\\$n_2,\theta^{(2)}$};
      \node[client,below right=14mm and 7mm of server] (c3) {Cliente 3\\$n_3,\theta^{(3)}$};
      \draw[feedback] (server)--(c1); \draw[feedback] (server)--(c2); \draw[feedback] (server)--(c3);
      \node[below=4mm of c2] {datos permanecen locales};
    \end{tikzpicture}
    \column{.44\textwidth}
    \small
    \[
      \theta_{t+1}=\sum_{k=1}^{K}\frac{n_k}{\sum_jn_j}\theta_t^{(k)}.
    \]
    \begin{enumerate}
      \item distribuir;
      \item ajustar localmente;
      \item enviar actualizaciones;
      \item agregar y repetir.
    \end{enumerate}
  \end{columns}
\end{frame}

% 34
\begin{frame}{Ejemplo federado: varias instituciones}
  \small
  Varios hospitales desean mejorar un clasificador sin centralizar historias clínicas.
  \begin{itemize}
    \item Cada institución conserva sus registros.
    \item Entrena localmente una copia del modelo.
    \item Comparte actualizaciones con un coordinador.
    \item El coordinador agrega y redistribuye el modelo.
  \end{itemize}
  \begin{alertblock}{Heterogeneidad}El resultado puede fallar si poblaciones, mediciones o tamaños difieren mucho entre instituciones.\end{alertblock}
\end{frame}

% 35
\begin{frame}{Federado no significa privado por sí solo}
  \small
  Las actualizaciones pueden revelar información y el servidor puede ser un punto de ataque.
  \begin{tabularx}{\textwidth}{>{\bfseries}p{.31\textwidth}X}\toprule
  Riesgo & Control posible\\\midrule
  exposición de actualizaciones & agregación segura y privacidad diferencial\\
  clientes maliciosos & autenticación y validación robusta\\
  poblaciones diferentes & evaluación global y por cliente\\
  costo de comunicación & selección de clientes y compresión\\
  clientes intermitentes & rondas tolerantes a disponibilidad\\\bottomrule
  \end{tabularx}
\end{frame}

% 36
\begin{frame}{Comparación final}
  \scriptsize
  \begin{tabularx}{\textwidth}{>{\bfseries}p{.15\textwidth}p{.22\textwidth}p{.29\textwidth}X}\toprule
  Método & Señal principal & Qué aprende & Salida típica\\\midrule
  SVM & etiquetas & frontera de margen & clase o puntuación\\
  Red neuronal & pérdida & transformaciones compuestas & predicción o representación\\
  Ensamble & errores etiquetados & combinación de modelos & predicción agregada\\
  RL & recompensas & política o valores & acción\\
  NLP / LLM & texto, etiquetas o tokens & representación y distribución & clase, extracción o texto\\
  Federado & actualizaciones locales & parámetros compartidos & modelo global\\\bottomrule
  \end{tabularx}
  \idea{Elegir exige alinear señal, salida, evaluación y condiciones de uso.}
\end{frame}

% 37
\begin{frame}{Síntesis}
  \small
  \begin{itemize}
    \item SVM selecciona una frontera mediante margen y ejemplos críticos.
    \item Las redes aprenden composiciones no lineales mediante gradientes.
    \item Los ensambles aprovechan diversidad entre predictores.
    \item RL aprende decisiones secuenciales a partir de recompensas.
    \item NLP convierte lenguaje en representaciones; los LLM generan token a token.
    \item El aprendizaje federado distribuye el ajuste, pero no garantiza privacidad.
  \end{itemize}
\end{frame}

% 38
\begin{frame}{Lecturas}
  \small
  \begin{block}{Modelos supervisados y distribuidos}
    \begin{itemize}
      \item Capítulo 12, §§12.2--12.4: SVM, ensambles y aprendizaje federado.
      \item Capítulo 15, §§15.1--15.2: redes neuronales y entrenamiento.
    \end{itemize}
  \end{block}
  \begin{block}{Decisiones y lenguaje}
    \begin{itemize}
      \item Capítulo 8, §§8.2 y 8.4: decisiones secuenciales y aprendizaje por refuerzo.
      \item Capítulo 16, §§16.1, 16.3 y 16.4: NLP, Transformer y modelos de lenguaje.
    \end{itemize}
  \end{block}
\end{frame}

\end{document}
